\section{Continuous Assessment 2 (CA-2) --- Set A}
\label{sec:ca2_set_a}

\LPUPaperHeader{Continuous Assessment (CA-2)}{A}{60 Minutes}{30 Marks}{CO4 (Ideation \& Selection), CO5 (Workflow Design \& Validation)}

\begin{center}
\sffamily\textbf{\color{AlertRed}Note: Attempt all six questions. Each question carries 5 marks.}
\end{center}

\questionitem{1}{
    Apply the SCAMPER creative ideation methodology to generate at least six innovative system concepts to redesign the physical campus library experience.
}{5}{CO4}{L3 Apply}

\questionitem{2}{
    Contrast traditional open-room brainstorming with the Brainwriting 6-3-5 protocol. Explain why Brainwriting produces higher conceptual variety and minimizes psychological bias.
}{5}{CO4}{L2 Understand}

\questionitem{3}{
    Explain the difference between a Wizard of Oz prototype and a Concierge MVP. Provide an illustrative software application example for each.
}{5}{CO5}{L2 Understand}

\questionitem{4}{
    Construct a Multi-Tier Service Blueprint for an online grocery home delivery platform. Diagram the Line of Interaction, Line of Visibility, and identify at least two Failure Nodes [!].
}{5}{CO5}{L4 Analyze}

\questionitem{5}{
    In a usability test of a new mobile banking interface across 10 participants, 8 participants successfully transferred funds without assistance, taking an average of 42 seconds. Calculate the Task Completion Rate (TCR) and evaluate the outcome against industry benchmarks.
}{5}{CO5}{L3 Apply}

\questionitem{6}{
    A usability evaluation yields the following 10 raw System Usability Scale (SUS) scores from a participant: $[5, 1, 4, 2, 5, 1, 4, 2, 5, 1]$. Calculate the overall composite SUS score and determine its usability grade.
}{5}{CO5}{L3 Apply}

\newpage
\subsection*{Solutions and Marking Rubrics --- CA-1 Set A}

\begin{solutionbox}[Solution to Question 1: SCAMPER Application]
\textbf{SCAMPER Redesign for Campus Library:}\\
$\bullet$ \textbf{Substitute (S):} Substitute static printed textbooks with scannable dynamic AR holograms and digital annotations.\\
$\bullet$ \textbf{Combine (C):} Combine silent study cubicles with sound-isolated podcast and video recording studios.\\
$\bullet$ \textbf{Adapt (A):} Adapt the Amazon fulfillment center logistics model by installing automated retrieval robots for archival books.\\
$\bullet$ \textbf{Modify / Magnify (M):} Magnify operating hours to 24/7/365 with automated biometric security pods.\\
$\bullet$ \textbf{Put to Other Use (P):} Repurpose underutilized basement book stacks into collaborative student incubation makerspaces.\\
$\bullet$ \textbf{Eliminate (E):} Eliminate manual checkout desks through RFID doorway arch sensors that auto-charge loans to student IDs.\\
$\bullet$ \textbf{Reverse / Rearrange (R):} Invert book acquisition: allow students to vote on library procurement budgets directly.
\begin{markingrubric}
$\bullet$ Six distinct SCAMPER categories correctly addressed: \textbf{3 Marks}\\
$\bullet$ Quality, feasibility, and originality of proposed concepts: \textbf{2 Marks}
\end{markingrubric}
\end{solutionbox}

\begin{solutionbox}[Solution to Question 2: Brainwriting 6-3-5 vs. Brainstorming]
\textbf{Pathologies of Traditional Brainstorming:}\\
1. \textit{Production Blocking:} Only one participant can speak at a time; others forget or abandon ideas while waiting.\\
2. \textit{Evaluation Apprehension:} Junior members fear looking foolish in front of senior leaders.\\
3. \textit{Social Loafing:} Disengaged participants hide behind talkative dominant individuals.\\
\textbf{The 6-3-5 Mechanism and Psychological Advantages:}\\
$\bullet$ 6 participants $\times$ 3 ideas $\times$ 6 rounds = 108 ideas in 30 minutes.\\
$\bullet$ Total silence eliminates vocal domination and interruptions.\\
$\bullet$ Round-robin passing forces systematic cross-pollination and evolutionary idea building without social hierarchy friction.
\begin{markingrubric}
$\bullet$ Identification of 3 brainstorming pathologies: \textbf{2 Marks}\\
$\bullet$ Explanation of 6-3-5 structure and psychological advantages: \textbf{3 Marks}
\end{markingrubric}
\end{solutionbox}

\begin{solutionbox}[Solution to Question 3: Wizard of Oz vs. Concierge MVP]
\begin{tabularx}{\linewidth}{|l|X|X|} \hline
\textbf{Dimension} & \textbf{Wizard of Oz Prototype} & \textbf{Concierge MVP} \\ \hline
\textbf{User Perception} & Front-end appears fully automated (AI/software-driven). & Service is explicitly disclosed as 100\% manual human white-glove support. \\ \hline
\textbf{Backstage Reality} & Humans secretly execute backend tasks manually behind the scenes. & Humans openly assist the customer directly in person or via call. \\ \hline
\textbf{Validation Goal} & Tests interface UX, trust, user journey, and perceived responsiveness. & Tests deep customer willingness to pay and customized service workflow. \\ \hline
\textbf{Example} & An automated legal contract redline app where human lawyers manually edit text within 10 minutes. & A personal shopping stylist who accompanies the user to clothing stores. \\ \hline
\end{tabularx}
\begin{markingrubric}
$\bullet$ Clear structural comparison across perception and backend: \textbf{3 Marks}\\
$\bullet$ Realistic, distinct software/business examples for both: \textbf{2 Marks}
\end{markingrubric}
\end{solutionbox}

\begin{solutionbox}[Solution to Question 4: Grocery Platform Service Blueprint]
\textbf{Service Blueprint Architectural Layers:}\\
$\bullet$ \textbf{Physical Evidence:} Mobile app UI catalog; order status tracker; SMS delivery pin; grocery paper bag.\\
$\bullet$ \textbf{Customer Actions:} Browses produce $\to$ adds to cart $\to$ selects delivery time slot $\to$ completes payment $\to$ receives groceries.\\
\textbf{--- Line of Interaction ---}\\
$\bullet$ \textbf{Frontstage Contact:} Delivery driver arrives at doorstep, hands over bag, confirms delivery pin.\\
\textbf{--- Line of Visibility ---}\\
$\bullet$ \textbf{Backstage Contact:} Warehouse picker gathers items with barcode scanner; packer bags items; dispatch manager routes driver.\\
\textbf{--- Line of Internal Interaction ---}\\
$\bullet$ \textbf{Support Processes:} Real-time warehouse inventory sync; credit card gateway settlement; automated route optimization algorithm.\\
$\bullet$ \textbf{Failure Nodes [!]:}\\
1. \textit{[! 1] Out-of-Stock Discrepancy:} In-app produce shows available, but warehouse bin is empty during picking.\\
2. \textit{[! 2] Payment Timeout Gateway Failure:} Bank server drops confirmation during session token handover.
\begin{markingrubric}
$\bullet$ Inclusion of all key blueprint tiers with lines of interaction/visibility: \textbf{3 Marks}\\
$\bullet$ Identification and analysis of 2 critical failure points: \textbf{2 Marks}
\end{markingrubric}
\end{solutionbox}

\begin{solutionbox}[Solution to Question 5: Task Completion Rate (TCR) Calculation]
\textbf{Formula:}
\[
\text{TCR} = \frac{\text{Successful Task Completions}}{\text{Total Task Attempts}} \times 100\% = \frac{8}{10} \times 100\% = \mathbf{80.0\%}
\]
\textbf{Benchmarking \& Evaluation:}\\
$\bullet$ In consumer mobile financial applications, an 80\% TCR represents an acceptable baseline, but falls below the 90\% threshold required for mission-critical core tasks (such as funds transfer).\\
$\bullet$ The 2 failed attempts (20\% drop-off) signal critical friction, likely caused by confusing beneficiary account number entry or OTP verification timeouts.
\begin{markingrubric}
$\bullet$ Correct formula and mathematical calculation: \textbf{2 Marks}\\
$\bullet$ Professional benchmarking evaluation and failure diagnosis: \textbf{3 Marks}
\end{markingrubric}
\end{solutionbox}

\begin{solutionbox}[Solution to Question 6: System Usability Scale (SUS) Score]
\textbf{Given Responses:} $R = [5, 1, 4, 2, 5, 1, 4, 2, 5, 1]$\\
Odd questions (1, 3, 5, 7, 9) contribute $(R_i - 1)$:\\
$\bullet$ $(5-1) + (4-1) + (5-1) + (4-1) + (5-1) = 4 + 3 + 4 + 3 + 4 = 18$\\
Even questions (2, 4, 6, 8, 10) contribute $(5 - R_i)$:\\
$\bullet$ $(5-1) + (5-2) + (5-1) + (5-2) + (5-1) = 4 + 3 + 4 + 3 + 4 = 18$\\
\textbf{Composite Calculation:}
\[
\text{SUS Score} = 2.5 \times (18 + 18) = 2.5 \times 36 = \mathbf{90.0}
\]
\textbf{Usability Grade Interpretation:}\\
A score of 90.0 places the system in the top 5\% of all benchmarked digital products (Grade A+, ``Best Imaginable'' usability), far exceeding the industry average of 68.0.
\begin{markingrubric}
$\bullet$ Correct formula application for odd and even items: \textbf{3 Marks}\\
$\bullet$ Accurate calculation resulting in 90.0: \textbf{1 Mark}\\
$\bullet$ Accurate interpretation against industry percentile benchmarks: \textbf{1 Mark}
\end{markingrubric}
\end{solutionbox}
